\documentclass[
    english,
	ruledheaders=section,%Ebene bis zu der die Überschriften mit Linien abgetrennt werden, vgl. DEMO-TUDaPub
	class=report,% Basisdokumentenklasse. Wählt die Korrespondierende KOMA-Script Klasse
	thesis={type=master},% Dokumententyp Thesis (projektseminar, proseminar, bachelor, master)
	accentcolor=2c, % Akzentfarbe FG rtm
	custommargins=false,
	marginpar=false,
	BCOR=10mm, % Bindekorrektur (etwas größere mittlere Ränder)
	parskip=half-, % Absatzkennzeichnung durch Abstand vgl. KOMA-Sript
	fontsize=11pt, % Basisschriftgröße laut Corporate Design ist mit 9pt häufig zu klein
  	twoside,
  	numbers=noendperiod,
  	fleqn,
  	captions=oneline, % Einzeilige Captions zentrieren
	captions=tableheading, % Für besseren Abstand bei TabellenÜBERschriften
  	IMRAD=false
]{tudapub}


% Der folgende Block ist nur bei pdfTeX auf Versionen vor April 2018 notwendig
\usepackage{iftex}
\ifPDFTeX
\usepackage[utf8]{inputenc}%kompatibilität mit TeX Versionen vor April 2018
\fi

%%%%%%%%%%%%%%%%%%%
%Sprachanpassung & Verbesserte Trennregeln
%%%%%%%%%%%%%%%%%%%
\usepackage[main=english, ngerman]{babel}
\usepackage[autostyle]{csquotes}% Anführungszeichen vereinfacht
\usepackage{microtype}


%%%%%%%%%%%%%%%%%%%
% Literaturverzeichnis
%%%%%%%%%%%%%%%%%%%
\catcode30=12  % Workaround for current biblatex problem (https://tex.stackexchange.com/questions/564990/error-after-miktex-reinstall-text-line-contains-an-invalid-character)
\usepackage[backend=biber,style=ieee,bibencoding=utf8]{biblatex}
\addbibresource{literature.bib}

% Additional Packages
\usepackage{acro} % Acronym management
\usepackage{amsmath, amssymb}
\usepackage{booktabs} % Professional table formatting
\usepackage{siunitx}
\usepackage{tikz}
\usepackage{graphicx} % for pdf, bitmapped graphics files
%\usepackage{cleveref} % Smart cross-referencing (load late)

% Acronym definitions (acro package)
\DeclareAcronym{mpc}{
  short = MPC ,
  long  = Model Predictive Control ,
}


% ======================================================
% Fix für Fehlermeldungen 'Undifined control sequence' und 'Missing number, treated as zero.' bei gestapelten Akzenten in der Mathe-Umgebung. (z.B. \dot{\hat{x}})
% Der Fehler tritt auf bei tuda-ci v3.10 (Feb. 2021) in Kombination mit dem amsmath package.
%
% BITTE PRÜFEN, OB DER FEHLER NOCH BESTEHT UND DEN WORKAROUND GGF. WIEDER ENTFERNEN
%
% workaround/solution by skillmon posted here
% https://tex.stackexchange.com/questions/482281/math-accent-symbol-over-parentheses-enclosing-accented-symbol-amsmath
\makeatletter
\protected\def\mathaccentV#1#2#3#4#5%
  {%
    \ifmmode
      \mathaccentV@do{#2}{#3}{#4}{#5}%
    \else
      \@xp\nonmatherr@\csname #1\endcsname
    \fi
  }
\def\mathaccentV@do#1#2#3#4%
  {%
    \global\let\macc@nucleus\@empty
    \mathaccent"\accentclass@#1#2#3{#4}\macc@nucleus
  }
\makeatother
% ======================================================

\begin{document}

\Metadata{	% Metadaten im erzeugten PDF/A
	title=Eine minimale LaTeX-Vorlage für schriftliche (Abschluss-)Arbeiten am IAT,
	author=Rudi Ratlos,
	keywords={TU Darmstadt \sep Corporate Design \sep LaTeX},
}

\title{Eine minimale \LaTeX-Vorlage für schriftliche (Abschluss-)Arbeiten am IAT}
\subtitle{\textmd{Electrical Engineering and Information Technology program}}
\author[R. Ratlos]{Rudi Ratlos}
\reviewer*[First Examiner,Supervisor]{Prof. Dr.-Ing. Rolf Findeisen \and NameGoesHere} % erst ab Version v3.06 (2020-10-20) des tuda-ci-Pakets möglich

\department{etit}

\addTitleBox{\includegraphics[width=\linewidth]{CCPS_logo}}

%\date{\today}
\submissiondate{1. April 2020}

\lowertitleback{
	Technische Universität Darmstadt\\
	Institut für Automatisierungstechnik und Mechatronik\\
	Fachgebiet Regelungstechnik und Cyber-Physische Systeme\\
	Prof. Dr.-Ing. Rolf Findeisen\\
}

\frontmatter
\maketitle

\clearpage
\setcounter{page}{1}
\section*{Task Description}

	Provided by supervisor

 \vspace{0.5cm}
\begin{tabular}{ll}
	Start date: 	& ------ \\
	End date: 		& ------ \\
	Date of the colloquium: 	& ------ \\
\end{tabular}
\clearpage

\affidavit[digital]

\section*{Abstract}
	% Abstract structure (increasing scope):
	% 1. Context (current state of the topic)
	% 2. Problem or open research points
	% 3. Solution
	% 4. Key takeaway

\section*{Zusammenfassung}
	Additional abstract in German.

\cleardoublepage

\tableofcontents
\cleardoublepage

\mainmatter

\chapter{Introduction}
\paragraph{Context}
\paragraph{Problem Statement and Open Research Points}
\paragraph{Proposed Solution}
\paragraph{Contributions}
\paragraph{Thesis Structure}

\chapter{Foundations}
\section{Whole-Brain Biophysical Modeling}

\subsection{Introduction \& Historical Context}
% Cite Jansen \& Rit (1995) and Lopes da Silva et al. (1974)
% Explain the neural mass modeling paradigm: mean-field approximation of cortical macro-columns

\subsection{Biological Architecture \& Microcircuit Setup}
% Describe the 3 interacting neuronal subpopulations:
%   1. Pyramidal cells (PY) - main projection neurons
%   2. Excitatory interneurons (eIN) - positive feedback loop
%   3. Inhibitory interneurons (iIN) - negative feedback loop
% Outline intrinsic connectivity constants (C_1, C_2, C_3, C_4)

\subsection{Mathematical Formulation of Single-Node Dynamics}

\subsubsection{Postsynaptic Impulse Responses \& Sigmoidal Activation}
% Linear pulse-to-wave operator: 2nd-order impulse response kernels h_e(t) and h_i(t)
% Static wave-to-pulse operator: instantaneous nonlinear sigmoid S(v)

\subsubsection{Six-State ODE Representation}
% State-space formulation with 6 coupled 1st-order ODEs per node (x_1 ... x_6)
% Physical interpretation of PSP states (x_1, x_2, x_3) and velocity states (x_4, x_5, x_6)
% Local node output signal: y_i(t) = x_{i,2}(t) - x_{i,3}(t) (net pyramidal depolarisation)

\subsubsection{Bifurcation Behavior \& Regional Excitability}
% Excitatory synaptic gain (A) as the primary bifurcation parameter
% Three-level excitability mapping (Yu et al. framework):
%   - Healthy background tissue (A = 3.25 mV): baseline alpha activity (~10 Hz)
%   - Propagation Zone / PZ (A = 3.4 mV): sub-critical, network-recruited state
%   - Epileptogenic Zone / EZ (A = 3.6 mV): autonomous seizure-like limit cycles
% Spectral dynamics, time constants (tau_e, tau_i), and bistability between limit cycles and fixed points

\subsection{Whole-Brain Connectome Network Coupling}

\subsubsection{Structural Connectivity \& Inter-Node Topology}
% Scaling from single column to N-region Brain Network Model (BNM)
% Connectome structural connectivity matrix (W_{ij}) from DTI / CoCoMac
% Global coupling scale factor (K)

\subsubsection{Axonal Conduction Delays}
% Inter-regional fiber tract lengths (d_{ij}) and conduction velocity (v)
% Calculation of conduction delays: \tau_{ij} = d_{ij} / v
% Integration of delayed presynaptic input S(y_j(t - \tau_{ij})) into excitatory PSP channel (x_{i,5})

\subsection{Electrophysiological Interpretation}
\paragraph{Primary and Secondary Dipole Currents}

\paragraph{EEG Observation Model}

\paragraph{Electrical Stimulation Model}
% Direct additive voltage perturbation inside the pyramidal sigmoid: S(x_{i,2} - x_{i,3} + U_{\text{tES},i})

% Reciprocity principle formulation and justification for having two different leadfield models for

\section{Data-Driven Dynamics Identification}

\subsection{Multilayer Perceptrons (MLP)}

\section{Model Predictive Control (MPC)}

\subsection{Fundamentals of Receding Horizon Control}

\subsection{Nonlinear ARX (NARX) Formulation}

\subsection{Mixed-Integer MPC for Discrete Neurostimulation}

\chapter{Methodology}
\section{Closed-Loop Neurostimulation Architecture}
\section{Surrogate Model Training and CasADi Integration}
\section{MPC Controller and Cost Function Design}

\chapter{Experiments and Evaluation}

\chapter{Conclusion and Discussion}
\section{Summary of Key Findings}
\section{Limitations and Critical Discussion}
\section{Future Directions}


\cite{Duden}

% =================================================================================
% Anhang
% =================================================================================
\appendix % Damit wird der Anhang begonnen. Die Kapitel werden ab jetzt mit Buchstaben nummeriert


% =================================================================================
% Literaturverzeichnis
% =================================================================================
\cleardoublepage        % Auf eine leere Seite einfügen
\phantomsection         % Für Aufnahme ins Inhaltsverzeichnis
\addcontentsline{toc}{chapter}{\bibname}  % In Inhaltsverzeichnis von
                                          % Dokument und pdf aufnehmen
\printbibliography
% =================================================================================


\end{document}
